\subsection{Diagramas de casos de uso}
\label{anexo1:sec:requisitos-funcionales-casos-de-uso}

\diagrama{
    nombre = {Casos de uso del sistema},
    archivo = {anexo1/diagramas/casosuso.pdf},
    etiqueta = {casos-uso},
    forzarPosicion = {sí pls},
}

\subsection{Definición de actores}
\label{anexo1:sec:requisitos-funcionales-actores}
En este subapartado se definen los distintos actores del sistema en forma de tabla, así como la
relacción entre ellos a través de un diagrama.

\diagrama{
    nombre = {Actores del sistema},
    archivo = {anexo1/diagramas/Actores.pdf},
    etiqueta = {actores},
}

\actor{
    id = {0001},
    nombre = {Usuario sin autenticar},
    descripcion = {
        Este actor representa a un usuario que no ha iniciado sesión en el sistema.
    },
    comentarios = {
        Este actor interactúa con el sistema unicamente para registrarse o iniciar sesión.
    }
}

\actor{
    id = {0002},
    nombre = {Usuario autenticado},
    descripcion = {
        Este actor representa a un usuario que ha iniciado sesión en el sistema.
    },
    comentarios = {
        Este usuario solo puede consultar sus datos personales, cerrar sesión y verificar su correo
        electrónico.
    }
}

\actor{
    id = {0003},
    nombre = {Usuario verificado},
    descripcion = {
        Este actor representa a un usuario que ha iniciado sesión en el sistema y ha verificado su 
        correo electrónico.
    },
    comentarios = {
        Este actor puede realizar todas las acciones disponibles en el sistema.
    }
}

\actor{
    id = {0004},
    nombre = {Sistema},
    descripcion = {
        Este actor representa al sistema en sí mismo.
    }
}


\subsection{Casos de uso del sistema}
\label{anexo1:sec:requisitos-funcionales-casos-de-uso-definicion}
\casouso{
    nombre = {Registrarse},
    descripcion = {%
        El sistema deberá comportarse tal como se describe en el siguiente caso de uso cuando el
        actor Usuario sin autenticar pulse el texto 'Regístrate' en la página principal o la página
        de inicio de sesión. 
    },
}{
    objetivosAsociados = {
        \item \nameref{anexo1:tab:obj0004}
        \item \nameref{anexo1:tab:obj0006}
    },
    requisitosAsociados = {
        \item \nameref{anexo1:tab:ri0001}
    },
    precondicion = {%
        El actor debe tener el rol de Usuario sin autenticar.
    },
    longsecuencia = {3},
    secuencia = {%
        & 1 & El sistema cambia el formulario de Iniciar Sesión por el de Registro \\
        & 2 & El usuario rellena el formulario y pulsa el botón de registrarse \\
        & 3 & El sistema valida la información introducida y, si es correcta, el usuario es creado \\
              % Mandar correo de verificación
              % Iniciar sesión
    },
    postcondicion = {%
        El actor Usuario sin autenticar pasa a tener el rol de Usuario autenticado.
    },
    longexcepciones = {1},
    excepciones = {%
        & 3 & Si el usuario ya existe, el sistema muestra un mensaje de error indicando que el
              correo electrónico ya está en uso \\
    },
}

\casouso{
    nombre = {Iniciar sesión},
    descripcion = {%
        El sistema deberá comportarse tal como se describe en el siguiente caso de uso cuando el
        actor Usuario sin autenticar intente iniciar sesión desde la página principal o la página
        de inicio de sesión, o durante la realización del caso de uso \nameref{anexo1:tab:cu1}
    },
}{}

\casouso{
    nombre = {Mandar correo de verificación},
    descripcion = {%
        El sistema deberá comportarse tal como se describe en el siguiente caso de uso abstracto
        durante la realización del caso de uso \nameref{anexo1:tab:cu1}
    },
}{}

\casouso{
    nombre = {Ver datos personales},
    descripcion = {%
        El sistema deberá comportarse tal como se describe en el siguiente caso de uso cuando el
        actor Usuario autenticado entre en la sección de datos personales ya sea desde el menú o
        justo después de iniciar sesión
    },
}{}

\casouso{
    nombre = {Eliminar cuenta},
    descripcion = {%
        El sistema deberá comportarse tal como se describe en el siguiente caso de uso cuando el
        actor Usuario autenticado pulse el botón de eliminar cuenta en el menú principal 
    },
}{}

\casouso{
    nombre = {Cerrar sesión},
    descripcion = {%
        El sistema deberá comportarse tal como se describe en el siguiente caso de uso cuando el
        actor Usuario autenticado pulse el botón de cerrar sesión en el menú principal o en la
        página principal
    },
}{}

\casouso{
    nombre = {Verificar correo electrónico},
    descripcion = {%
        El sistema deberá comportarse tal como se describe en el siguiente caso de uso cuando el
        actor Usuario autenticado pulse el botón de verificar correo electrónico en el menú
        principal
    },
}{}

\casouso{
    nombre = {Modificar datos personales},
    descripcion = {%
        El sistema deberá comportarse tal como se describe en el siguiente caso de uso cuando el
        Usuario verificado modifique alguno de sus datos personales en el menú principal
    },
}{}

\casouso{
    nombre = {Añadir forma física},
    descripcion = {%
        El sistema deberá comportarse tal como se describe en el siguiente caso de uso cuando el
        Usuario verificado pulse el botón de añadir forma física en el menú principal
    },
}{}

\casouso{
    nombre = {Consultar estadísticas},
    descripcion = {%
        El sistema deberá comportarse tal como se describe en el siguiente caso de uso cuando el
        actor Usuario verificado pulse el botón de consultar estadísticas en el menú principal
    },
}{}

\casouso{
    nombre = {Filtrar estadísticas},
    descripcion = {%
        El sistema deberá comportarse tal como se describe en el siguiente caso de uso cuando el
        actor Usuario verifcado pulse alguno de los filtros temporales disponibles en el menú de
        estadísticas
    },
}{}

\casouso{
    nombre = {Abrir el juego},
    descripcion = {%
        El sistema deberá comportarse tal como se describe en el siguiente caso de uso cuando el
        actor Usuario verificado pulse el botón de jugar en el menú principal
    },
}{}

\casouso{
    nombre = {Activar pantalla completa},
    descripcion = {%
        El sistema deberá comportarse tal como se describe en el siguiente caso de uso cuando el
        actor Usuario verificado pulse el botón de activar pantalla completa en el menú de juego.
    },
}{}

\casouso{
    nombre = {Calibrar cámara},
    descripcion = {%
        El sistema deberá comportarse tal como se describe en el siguiente caso de uso abstracto una
        vez se inicie el caso de uso \nameref{anexo1:tab:cu12}
    },
}{}

\casouso{
    nombre = {Escoger minijuego},
    descripcion = {%
        El sistema deberá comportarse tal como se describe en el siguiente caso de uso una vez termine
        el caso de uso \nameref{anexo1:tab:cu13} o \nameref{anexo1:tab:cu20}
    },
}{}

\casouso{
    nombre = {Escoger dificultad},
    descripcion = {%
        El sistema deberá comportarse tal como se describe en el siguiente caso de uso una vez termine
        el caso de uso \nameref{anexo1:tab:cu14} o \nameref{anexo1:tab:cu19}
    },
}{}

\casouso{
    nombre = {Jugar minijuego},
    descripcion = {%
        El sistema deberá comportarse tal como se describe en el siguiente caso de uso una vez termine
        el caso de uso \nameref{anexo1:tab:cu15}
    },
}{}

\casouso{
    nombre = {Ver puntuaciones},
    descripcion = {%
        El sistema deberá comportarse tal como se describe en el siguiente caso de uso una vez termine
        el minijuego escogido por el usuario, es decir, al terminar el caso de uso \nameref{anexo1:tab:cu16}
    },
}{}

\casouso{
    nombre = {Volver a jugar},
    descripcion = {%
        El sistema deberá comportarse tal como se describe en el siguiente caso de uso cuando el actor
        Usuario verificado pulse el botón de volver a jugar en el menú de puntuaciones del minijuego.
    },
}{}

\casouso{
    nombre = {Volver a la selección de minijuego},
    descripcion = {%
        El sistema deberá comportarse tal como se describe en el siguiente caso de uso cuando el actor
        Usuario verificado pulse el botón de volver al menú principal en el menú de puntuaciones del
        minijuego.
    },
}{}

